\section*{РОЗДІЛ 5. ПРАКТИЧНА АПРОБАЦІЯ: СИСТЕМА RESEARCHIQ}

\subsection*{5.1. Архітектура рішення}

\subsubsection*{5.1.1. Концептуальна модель}

Система \textbf{ResearchIQ} (раніше Research Tracker Pro) реалізує запропоновану концептуальну модель інтеграції GenAI для \textit{research information management} та наукового адміністрування. В основі лежить принцип \textbf{агентної оркестрації}: LLM використовується для інтерпретації запитів природною мовою, планування послідовності інструментальних дій та пояснення результатів, тоді як збирання, нормалізація, дедуплікація та злиття даних реалізуються переважно детермінованими процедурами (див. розд.~4).

\textbf{Трирівнева архітектура:}

\begin{longtable}{@{}p{0.17\textwidth}p{0.42\textwidth}p{0.33\textwidth}@{}}
\toprule
\textbf{Рівень} & \textbf{Компоненти} & \textbf{Функції} \\
\midrule
\textbf{Data Layer} & ORCID API v3, OpenAlex REST API, (Scopus*, WoS*) & Отримання наукометричних даних з зовнішніх джерел \\
\textbf{Agent Layer} & LLM + MCP Function Calling & Оркестрація викликів tools, інтеграція результатів, пояснення \& діалогова взаємодія \\
\textbf{Interface Layer} & React Dashboard, Chat Interface, Report Generator & Взаємодія з користувачем, візуалізація, експорт \\
\bottomrule
\end{longtable}

\textit{Примітка: Scopus та WoS інтегруються опціонально при наявності інституційної підписки.}

\subsubsection*{5.1.2. Технологічний стек}

\textbf{Frontend:}
\begin{itemize}[leftmargin=*,itemsep=0.2em]
\item \textbf{React 18} з TypeScript для типобезпечного коду;
\item \textbf{Tailwind CSS} для responsive дизайну;
\item \textbf{Recharts} для візуалізації даних;
\item \textbf{React Query} для кешування та синхронізації стану.
\end{itemize}

\textbf{Backend / AI:}
\begin{itemize}[leftmargin=*,itemsep=0.2em]
\item LLM з нативною підтримкою Function Calling;
\item streaming responses для взаємодії в режимі реального часу.
\end{itemize}

\textbf{Інтеграції:}
\begin{itemize}[leftmargin=*,itemsep=0.2em]
\item \textbf{ORCID Public API v3} --- отримання профілів та works дослідників;
\item \textbf{OpenAlex REST API} --- розширені метадані, цитування, концепти;
\item \textbf{Crossref API} --- верифікація DOI та метаданих (коли доступно).
\end{itemize}

\subsubsection*{5.1.3. Механізм MCP Function Calling}

Ключовим інженерним механізмом є імплементація \textbf{Function Calling} у зв'язці з \textbf{Model Context Protocol (MCP)}: це дозволяє LLM викликати зовнішні інструменти через формалізовані схеми та уніфікований контракт, а також інтегрувати їхні відповіді у кінцевий результат.

\textit{Приклад декларації tools (TypeScript) наведено у Додатку B (лістинг~\ref{lst:b-tools}).}

\textbf{Логіка оркестрації:} коли користувач вводить запит природною мовою, агент (1) аналізує намір; (2) визначає, які tools потрібні; (3) формує виклики з параметрами; (4) отримує результати та інтегрує їх у відповідь; (5) за потреби ініціює контур \textit{human-in-the-loop} для підтвердження критичних дій (напр., злиття записів).

\subsubsection*{5.1.4. Алгоритм Golden Record в дії}

\textbf{Крок 1: Збір з ORCID}

\textit{Логи виконання кроків 1--3 та підсумковий результат наведено у Додатку B (лістинги~\ref{lst:b-orcid-log}--\ref{lst:b-result-log}).}

\subsection*{5.2. Дизайн оцінювання та попередні спостереження (pilot)}

\subsubsection*{5.2.1. Контекст пілоту та вибірка}

У межах пілотного впровадження система застосовувалася для консолідації бібліографічних даних для \textbf{15} науково-педагогічних працівників (\textbf{3} кафедри одного ЗВО) за період \textbf{2020--2024} рр. Вибірка включала різні наукові напрямки та різну інтенсивність публікацій (від поодиноких публікацій до профілів із десятками/сотнями записів).

\subsubsection*{5.2.2. Метрики та операціональні визначення (evaluation protocol)}

Нижче наведено мінімальний протокол оцінювання, який робить твердження про ефективність відтворюваними. \textbf{Зауваження:} у поточній версії рукопису протокол зафіксовано, але \textbf{повні репліковані виміри (N, дисперсія/CI) та gold-розмітка ще не наведені} (\textit{Evidence gap}).

\textbf{Одиниця спостереження.} Операція (task) типу: (i) формування списку публікацій для одного дослідника; (ii) формування річного звіту для кафедри; (iii) дедуплікація/злиття набору записів.

\textbf{Час виконання.} Для кожного task фіксуються:
\begin{itemize}[leftmargin=*,itemsep=0.2em]
\item момент старту/завершення (timestamp) та загальна тривалість;
\item конфігурація середовища (мережа, кеш увімкн./вимкн.);
\item кількість зовнішніх запитів (ORCID/OpenAlex/Crossref) та кількість помилок/повторів (retries).
\end{itemize}

\textbf{Дублікати (duplicate rate).} Нехай $R$ --- множина отриманих записів, $C$ --- кластери, утворені алгоритмом Golden Record. Дублікатами вважаються записи, що належать до одного кластера, але представляють одну й ту саму реальну сутність (публікацію). Для відтворюваного вимірювання потрібна \textbf{gold-розмітка} на підвибірці:
\begin{itemize}[leftmargin=*,itemsep=0.2em]
\item \textbf{Pairwise metrics} для зіставлення: precision/recall/F1;
\item \textbf{Clustering metrics} (напр., pairwise-F1 або B-cubed F1) для оцінки кластерів;
\item \textbf{Error analysis}: типові false merge / missed merge кейси.
\end{itemize}

\textbf{Повнота/точність полів.} Для атрибутів (DOI, рік, назва, venue) оцінюється частка правильних значень щодо еталону; еталоном може бути (a) ручна верифікація, (b) пріоритетне зовнішнє джерело (напр., Crossref для DOI) за зафіксованою політикою довіри (див. розд.~4).

\textbf{Бейзлайни та абляції (план).} Для ізоляції внесків плануються порівняння з базовими правилами (DOI-only, Title+Year) та/або ER-tool, а також абляції: (i) детермінований оркестратор vs MCP-агент; (ii) LLM-everywhere vs selective-LLM.

\subsubsection*{5.2.3. Попередні спостереження (без строгих claims)}

У пілоті було отримано практичні свідчення зменшення ручної роботи під час консолідації та форматування звітності, зокрема за рахунок: (i) автоматизованого збирання з кількох джерел, (ii) детермінованої нормалізації та дедуплікації, (iii) уніфікованих шаблонів експорту.

\textbf{Важливо:} конкретні числові значення виграшів у часі та якості даних залежать від інструментування, кешування, ступеня заповненості ORCID-профілів і політик зіставлення; у поточному рукописі вони розглядаються як \textbf{гіпотези, що потребують реплікованого підтвердження} (\textit{Evidence gap}).

\subsection*{5.3. Якісні результати та сценарії використання}

\subsubsection*{5.3.1. Перехід до проактивності}

Якісний ефект від використання ResearchIQ пов'язаний зі зміною режиму роботи з дослідницькими метаданими: від епізодичного, реактивного формування звітів до \textit{постійного} оновлення, моніторингу та запитно-орієнтованої аналітики.

\subsubsection*{5.3.2. Аналітичні можливості (illustrative examples)}

Аналітичний модуль орієнтовано на відповіді на запити природною мовою та подальшу верифікацію результатів у вигляді агрегованих таблиць і графіків.

\textbf{Примітка: наведені нижче діалоги є демонстраційними (synthetic) прикладами для ілюстрації класів запитів і формату відповіді; вони не є прямими вимірами на реальних даних пілоту} (\textit{Evidence gap щодо конкретних значень}).

\textbf{Порівняльний аналіз.}
\begin{quote}
\textit{Запит:} «Хто з нашої кафедри має найвищий H-index та як він змінився за останні 3 роки?»

\textit{Відповідь агента (приклад формату):} «…»
\end{quote}

\textbf{Тематичний аналіз.}
\begin{quote}
\textit{Запит:} «Які теми найчастіше зустрічаються в публікаціях нашого факультету?»

\textit{Відповідь агента (приклад формату):} «…»
\end{quote}

\textbf{Прогнозування / аномалії (статус прототипу).}
\begin{quote}
\textit{Запит:} «Які тренди публікаційної активності спостерігаються та чи є нетипові відхилення?»

\textit{Відповідь агента (приклад формату):} «…»
\end{quote}

\textit{Зауваження:} модулі прогнозування та виявлення аномалій у поточній роботі описуються як \textbf{прототипні можливості}; їхня якість потребує окремого протоколу оцінювання (помилки I/II роду, користь для користувачів), який виходить за межі цього пілоту (\textit{Evidence gap}).

\subsubsection*{5.3.3. Гнучкість звітності}

Уніфікована генерація звітів із Golden Record у кількох вихідних форматах (DOCX/XLSX/PDF) зменшує витрати на повторне форматування та знижує ризик помилок перенесення даних між системами. Додатково підтримуються кастомні інституційні шаблони, які заповнюються автоматично на основі агрегованих даних.

\subsubsection*{5.3.4. Human-in-the-Loop (HITL) верифікація}

Для процесів, що мають адміністративні наслідки, принципово важливо зберігати перевірюваність і контроль. Тому ResearchIQ реалізує контур \textbf{human-in-the-loop}: перед формуванням фінального звіту користувач отримує попередній перегляд; записи з низькою впевненістю маркуються для ручної перевірки; будь-яке автоматичне рішення може бути перевизначене; ключові дії фіксуються в аудит-лозі.

\subsection*{5.4. Обмеження, загрози валідності та lessons learned}

\subsubsection*{5.4.1. Виявлені обмеження}

Апробація виявила низку обмежень, зумовлених як якістю вхідних даних, так і характеристиками зовнішніх сервісів. Зокрема: (i) залежність від актуальності та повноти інтегрованих API; (ii) обмеження на частоту запитів (rate limits), що впливають на масові операції; (iii) ефект «холодного старту» для дослідників без ORCID або з неповними профілями; (iv) мовні та покриттєві обмеження (відмінності в індексації між джерелами).

\subsubsection*{5.4.2. Threats to validity (DSR framing)}

\textbf{Construct validity.} Без формалізованих визначень метрик (duplicate rate, completeness, accuracy) та gold-розмітки існує ризик неоднозначної інтерпретації «якості даних». \textit{Мітигація:} у розд.~5.2.2 зафіксовано протокол і необхідні метрики; потрібна його емпірична реалізація.

\textbf{Internal validity.} Оцінки часу чутливі до кешування, мережі, rate limits, а також до відмінностей у ручному «до» процесі. \textit{Мітигація:} логування та повторювані виміри за контрольованих умов.

\textbf{External validity.} Пілот на одному ЗВО може не узагальнюватися на інші установи з іншими дисциплінарними профілями та практиками ведення ORCID. \textit{Мітигація:} мультисайт-оцінювання та стратифікована вибірка.

\subsubsection*{5.4.3. Ключові уроки}

Насамперед, \textbf{якість вхідних даних} (заповненість ORCID, наявність DOI, консистентність назв) залишається ключовим чинником результативності. Ефективним є підхід поступового впровадження через пілотну групу з подальшим масштабуванням після стабілізації правил зіставлення та процедур верифікації. Важливою є комунікація з користувачами: система не замінює експертів, а зменшує частку рутинної роботи, підвищуючи можливості стратегічного планування.
